\noindent\textbf{Scope and editorial basis.}
This chapter covers all 27 MCCME records 462--488 in manuscript source order, with
dated assessments as of 10 November 2026. The Russian transcription is the textual
authority; the scans arbitrate doubtful formulas and attribution. The supplied
79-page envelope includes three boundary pages: P000840 is already represented by
1978-460.12--1978-460.20, and P000917--P000918 belong to the session dated
22 September 1980. Neither is counted again as a 1979--80 occurrence. The present
year's content occupies P000841--P000916. Record 481 is retained as a catalog
heading. The critical date of record 473 is 29 October 1979, against the catalog's
29 September; on 19 November, Gelfand's record 477 precedes Beilinson's record 476
in the manuscript.

The occurrence index contains 252 distinct mathematical items, 10
problem threads, and 104 bibliography entries. Item boundaries distinguish
definitions, constructions, mathematical claims, explicit questions, and identifiable
interventions. A class records historical fate; literal correctness and strength of
publication match are separate assessments.
\begin{center}
\begin{tabular}{lrrrrrr}
\toprule
Class & A & B & C & D & E & F\\
\midrule
Items & 12 & 7 & 3 & 0 & 203 & 27\\
\bottomrule
\end{tabular}
\end{center}

\noindent\textbf{The main lines.}
Much of the autumn is an extended mathematical encounter with quantum field theory:
Kallosh's reports on gravity and supergravity; Batalin's careful progression through
constraints, gauge slices, ghosts and BRST identities; and Kobzarev's account of
gauge fields, confinement and particle masses. Formal functional integrals are kept
separate from convergent finite-dimensional formulas. The early BRST material is
not retrospectively identified wholesale with the later BV formalism.

The surrounding mathematics is unusually broad. Zakharov treats resonance and
integrability in September and long-wave hydrodynamics in February. Gindikin's
twistor geometry and Berezin's supermathematics lead into projective bundles and
exterior algebras. Shabat's double-translation surfaces require the missing
nondegeneracy hypotheses in the classical Jacobian theorem. H\o egh-Krohn's energy
representations retain a delicate two-dimensional question. In November, Gelfand's
dilogarithm list is followed by Beilinson's specific regulator construction; the
latter has a close 1980 publication match. January brings Migdal's loop equations,
Zelevinsky's Kazhdan--Lusztig material and MacPherson's topology. Olshanskii's
representations and MacPherson's combinatorial characteristic classes distinguish
pre-talk results from genuinely subsequent developments. Their precise source
locators and chronology are given in the corresponding entries.

\noindent\textbf{What remains and what is developed here.}
The chapter does not turn a broad research programme into a definite new conjecture.
Continuum confinement is linked to the older static-potential problem; a closed
large-rank equation is not itself an area-law proof. The general two-dimensional
energy-representation question is distinguished from restricted cases and from a
published proof claim with a documented gap. Negative literature searches are reported as retrieval
limits, not proof that an idea was abandoned.

Berezin's two explicit matrix questions support a bounded mathematical treatment.
The normal-form question has an algorithmic answer through a finite-dimensional
matrix-algebra action; an explicit matrix-pair reduction explains the limits of a
simple Jordan analogy. For the even commutant of an arbitrary odd two-by-two matrix
over four Grassmann generators, the development gives a complete coefficient
parametrization, all dimension strata, and full proofs. Its exact computations were
implemented in two independent ways and agree in all 36 cases. These are proofs
and checks supplied here, with no priority claim. The literal finite commutant task
is elementary linear algebra, classified E; no open problem is manufactured to
justify its explanatory development. A separate compactly supported SU(2) example
checks the documented flaw in the energy-representation argument, without claiming
to resolve irreducibility.

\noindent\textbf{Continuity and corrections.}
The earlier occurrence index supplies links to exterior-module descriptions of
projective bundles, combinatorial Pontryagin classes, continuum confinement, and
the preceding year's integrable-system discussions. A common topic is not enough
to merge two problems with different quantifiers. The 106 corrections distinguish
notebook errors faithfully transcribed, unresolved readings, and provenance
discrepancies. In particular the adjacent-year
spillover, the October date, regulator signs, normalized Wilson traces, and the
Pontryagin construction's low-dimensional cases must remain visible to a reader.
The report gives the counting rules, evidence limits, search record and website
presentation recommendations.
